%! Absolute Value Equations \& Inequalities — Unit 1, Lesson 1.2 — Experience & Formalize (Activity · QuickNotes · Application · Check Your Understanding)
% Math Medic sizing: 12pt body, OPEN white space after every question (no writing lines).
% \answerspace{H}{} reserves exactly H in the blank; the key passes the answer as the second
% argument, so the two files paginate identically by construction.
% Page budget (user, 2026-08-20): activity 2pp max · QuickNotes 1/2pp · Application 1/2-1pp ·
% Check Your Understanding 1-2pp (practice, NOT scored).
\documentclass[12pt]{article}
\usepackage{algebra2-article}
\usepackage{algebra2-boxes}
\newcommand{\answerspace}[2]{\par\nopagebreak\noindent\begin{minipage}[t][#1][t]{\linewidth}%
  \color{keyred}\bfseries #2\end{minipage}\par}
% Plain number line: {xmin}{xmax}, ticks labeled every unit.
\newcommand{\numline}[2]{%
  \draw[->, semithick, charcoal] (#1-0.5,0)--(#2+0.5,0);
  \foreach \x in {#1,...,#2}{\draw[charcoal] (\x,0.12)--(\x,-0.12) node[below, font=\scriptsize]{$\x$};}}
% The trail: mile markers 0..20, ticks every mile, labels every two, tower at mile 12.
\newcommand{\trail}{%
  \draw[->, semithick, charcoal] (-0.7,0)--(21.2,0) node[right, font=\scriptsize]{mile};
  \foreach \x in {0,...,20}{\draw[charcoal] (\x,0.26)--(\x,-0.26);}
  \foreach \x in {0,2,...,20}{\node[below, font=\scriptsize, charcoal] at (\x,-0.30){$\x$};}
  \fill[forest] (12,0) circle (3pt);
  \node[forest, font=\scriptsize, above] at (12,0.34){tower};}

\begin{document}

\pageheader{Unit 1, Lesson 1.2}{Experience \& Formalize: In Range}

\vspace{0.06in}

\begin{headlinebox}{forestbg}
\textbf{In Range.} A straight stretch of parkway is marked off in miles, $0$ through $20$. A ranger's
radio tower stands at \textbf{mile 12}, and a handheld radio can reach it only from
\textbf{within 5 miles}. With your group, work through the two problems below using what you already
know --- number lines, distance, and solving --- and be ready to explain \emph{how you know} each
answer from the \textbf{picture}.
\end{headlinebox}

\vspace{0.04in}

\begin{scenariobox}[1. Exactly Five Miles Out]{forest}
\begin{enumerate}[label=\textbf{\alph*.}, leftmargin=1.8em, itemsep=8pt]
  \item Two hikers are standing \emph{exactly} $5$ miles from the tower. Mark both of them on the
        trail, then give their mile markers.
    \begin{center}
    \begin{tikzpicture}[x=0.62cm, y=0.5cm]
      \trail
    \end{tikzpicture}
    \end{center}
    Mile markers: \blank{2.0cm} and \blank{2.0cm}
  \item Why are there \emph{two} answers here and not one? Answer from the picture.
    \answerspace{1.8cm}{}
  \item A hiker stands at mile $x$. Write an expression, using absolute-value bars, for how far that
        hiker is from the tower. \quad Distance $=$ \blank{3.2cm}
  \item So ``exactly $5$ miles from the tower'' is the equation \blank{3.2cm} $= 5$. \\[4pt]
        Substitute each of your two markers from part (a) into it and show that both really work.
    \answerspace{2.0cm}{}
\end{enumerate}
\end{scenariobox}

\boxguard[15]
\begin{scenariobox}[2. In Range, Out of Range]{forest}
Remember: the radio reaches the tower only from within $5$ miles of mile $12$.

\begin{enumerate}[label=\textbf{\alph*.}, leftmargin=1.8em, itemsep=8pt]
  \item \textbf{Shade} every spot on the trail where a hiker's radio \emph{can} reach the tower.
    \begin{center}
    \begin{tikzpicture}[x=0.62cm, y=0.5cm]
      \trail
    \end{tikzpicture}
    \end{center}
    The shading runs from mile \blank{1.6cm} to mile \blank{1.6cm}, which as an inequality in $x$
    --- no bars --- is \blank{4.4cm}
  \item Now write ``within $5$ miles of the tower'' \emph{with} the bars, using your expression from
        1c. \quad \blank{3.4cm} $\le 5$
  \item \textbf{Shade} every spot where the radio \emph{cannot} reach the tower, then write those
        markers as an inequality (or inequalities) in $x$.
    \begin{center}
    \begin{tikzpicture}[x=0.62cm, y=0.5cm]
      \trail
    \end{tikzpicture}
    \end{center}
    Out of reach when \blank{6.0cm}
  \item Compare your two answers. The ``can reach'' markers are \emph{one} connected stretch of
        trail, but the ``cannot reach'' markers took \emph{two} separate pieces. Explain why, using
        the picture.
    \answerspace{2.4cm}{}
  \item A hiker at mile $3$ tries to radio in. Use your expression from 1c to decide whether the
        call gets through, and say how the number you computed tells you.
    \answerspace{2.0cm}{}
  \item The parkway only exists between mile $0$ and mile $20$. Does that change your answer in
        part (c)? Say exactly what has to be trimmed, and why.
    \answerspace{1.8cm}{}
\end{enumerate}
\end{scenariobox}


% ── QuickNotes — target: half a page ─────────────────────────────────────────
\boxguard[14]
\begin{tcolorbox}[enhanced, colback=sky, colframe=navy, boxrule=0.8pt, arc=2mm,
    left=3mm, right=3mm, top=1.5mm, bottom=1.5mm, breakable,
    fonttitle=\bfseries\small\color{white}, title={QuickNotes: Absolute Value Equations \& Inequalities},
    attach boxed title to top left={yshift=-2mm, xshift=4mm},
    boxed title style={colback=navy, colframe=navy, arc=1.5mm}]
\small
\begin{minipage}[t]{0.31\linewidth}
\vspace{2pt}\centering
\begin{tikzpicture}[x=0.40cm, y=0.42cm]
  \numline{-2}{8}
  \draw[very thick, greenacc] (1,0)--(5,0);
  \draw[very thick, redacc] (-2.5,0)--(1,0);
  \draw[very thick, redacc] (5,0)--(8.5,0);
  \fill[charcoal] (1,0) circle (2.6pt);
  \fill[charcoal] (5,0) circle (2.6pt);
  \node[greenacc, font=\tiny, above] at (3,0.5){$|x-3|<2$};
  \node[redacc, font=\tiny, above] at (-0.7,0.2){$>2$};
  \node[redacc, font=\tiny, above] at (6.7,0.2){$>2$};
\end{tikzpicture}\\[1pt]
{\footnotesize Distance from $3$, with $c=2$}
\end{minipage}\hfill
\begin{minipage}[t]{0.65\linewidth}
\vspace{0pt}
\begin{itemize}[leftmargin=1.1em, itemsep=4pt, topsep=0pt]
  \item $|x-h|$ is the \textbf{distance} between $x$ and \blank{0.9cm}; a distance is never
        \blank{1.7cm}.
  \item \textbf{Two-case rule.} $|X| = c$ with $c>0$: \; $X = $ \blank{1.0cm} or $X = $
        \blank{1.0cm}. \\
        $c=0$: \blank{1.0cm} solution. \quad $c<0$: \blank{1.0cm} solution.
  \item \textbf{Isolate first.} $3|x+1|-2=10$ splits only after $|x+1| = $ \blank{0.9cm}.
\end{itemize}
\end{minipage}

\vspace{4pt}
\begin{itemize}[leftmargin=1.1em, itemsep=4pt, topsep=0pt]
  \item \textbf{``Less th\emph{AND}''} $(<,\le)$: $|X|<c$ says $X$ is \emph{closer} than $c$ to
        zero, so \; $\blank{0.9cm} < X < \blank{0.9cm}$ --- one interval (\emph{between}).
  \item \textbf{``Great\emph{OR}''} $(>,\ge)$: $|X|>c$ says $X$ is \emph{farther} than $c$, so \;
        $X < \blank{0.9cm}$ \, or \, $X > \blank{0.9cm}$ --- two rays (\emph{outside}).
  \item \textbf{Three notations, one set.} For $1 \le x \le 5$: \; set $\{x \mid \blank{2.2cm}\}$,
        \; interval \blank{1.8cm}, \; and a shaded number line. Brackets and \blank{1.4cm} dots
        include an endpoint; parentheses and \blank{1.4cm} dots exclude it ($\pm\infty$: always a
        parenthesis).
\end{itemize}
\end{tcolorbox}

% ── Application — worked together, target: half to one page ──────────────────
\boxguard[14]
\begin{notesbox}{Application: Roasting Tolerance}
We will do this one together. A roaster fills bags labeled \textbf{340 grams}. Quality control
\textbf{rejects} any bag whose weight is \textbf{more than 8 grams off} the label.

\begin{enumerate}[leftmargin=1.7em, itemsep=7pt]
  \item Let $w$ be a bag's weight in grams. Write an absolute value inequality that says the bag is
        \emph{acceptable}. \quad \blank{4.2cm}
  \item Solve it.
    \begin{work}
             |w - 340| &\le 8 \\
        -8 \le w - 340 &\le 8 \\
             332 \le w &\le 348
    \end{work}
    Lightest acceptable bag: \blank{2.0cm} g \qquad Heaviest: \blank{2.0cm} g
  \item A bag comes off the line at $331$ g. Accepted or rejected --- and how does your answer to
        item~2 tell you, without re-solving anything?
    \answerspace{1.6cm}{}
  \item The roaster tightens the rule to ``more than $5$ grams off.'' Does the acceptable range get
        \emph{wider} or \emph{narrower}? \blank{2.4cm} \; What are the new endpoints?
        \blank{3.4cm}
\end{enumerate}
\end{notesbox}

% ── Check Your Understanding — practice, NOT scored ──────────────────────────
% This is the lesson's in-class practice set (graded homework is in DeltaMath): the six items
% below carry the whole A2.EI.1 spread (equation, "and", "or", special cases,
% tolerance model, and the formative MC).
\boxguard[14]
\begin{notesbox}{Check Your Understanding \quad {\normalfont\itshape (practice --- not scored)}}
Six exercises --- the lesson's practice.  Use these if you are having trouble, and bring to study
hall or office hours to get some help. Isolate the bars \emph{before} you split.

\begin{enumerate}[leftmargin=1.7em, itemsep=7pt]
  \item \textbf{Isolate, split, verify.} Solve $3|x+1| - 2 = 10$, then check one solution in the
        original equation.
    \begin{work}
      3|x+1| - 2 &= 10 \\
          3|x+1| &= 12 \\
           |x+1| &= 4 \\
             x+1 &= 4 \quad\text{or}\quad x+1 = -4 \\
               x &= 3 \quad\text{or}\quad x = -5
    \end{work}
    Solutions: $x = $ \blank{3.0cm} \qquad Check:
    \answerspace{1.0cm}{}

  \item \textbf{``Less th\emph{AND}.''} Solve $|2x - 1| \le 5$, then write the set three ways.
    \begin{work}
              |2x-1| &\le 5 \\
        -5 \le 2x-1 &\le 5 \\
           -4 \le 2x &\le 6 \\
            -2 \le x &\le 3
    \end{work}
    Set: \blank{3.4cm} \quad Interval: \blank{2.4cm} \quad
    \begin{tikzpicture}[x=0.44cm, y=0.44cm, baseline=-2pt]
      \numline{-5}{5}
    \end{tikzpicture}

  \boxguard[10]
  \item \textbf{``Great\emph{OR}.''} Solve $|3x - 6| \ge 9$, then write the set three ways.
    \begin{work}
             |3x-6| &\ge 9 \\
             3x - 6 &\le -9 \quad\text{or}\quad 3x - 6 \ge 9 \\
                 3x &\le -3 \quad\text{or}\quad 3x \ge 15 \\
                  x &\le -1 \quad\text{or}\quad x \ge 5
    \end{work}
    Set: \blank{4.0cm} \quad Interval: \blank{3.4cm} \\[4pt]
    \hspace*{1em}
    \begin{tikzpicture}[x=0.44cm, y=0.44cm, baseline=-2pt]
      \numline{-4}{8}
    \end{tikzpicture}

  \item \textbf{The odd ones out.} Answer each from the \emph{distance} meaning --- not a memorized
        rule. \\[3pt]
        (a) $|x - 4| = -3$ has \blank{3.4cm} \qquad (b) $|x - 4| = 0$ has \blank{3.4cm}
    \answerspace{1.2cm}{}
        (c) For which values of $k$ does $|x - 3| < k$ have \emph{no} solution? \blank{2.6cm}
    \answerspace{1.0cm}{}

  \boxguard[12]
  \item \textbf{Model it.} A thermostat holds a room within $2^\circ$F of its $68^\circ$F setting.
    \begin{enumerate}[label=(\alph*), leftmargin=1.9em, itemsep=5pt]
      \item Write an absolute value inequality for an acceptable temperature $T$. \blank{4.0cm}
      \item Solve it.
        \begin{work}
                |T - 68| &\le 2 \\
          -2 \le T - 68 &\le 2 \\
                66 \le T &\le 70
        \end{work}
        Coolest \blank{2.0cm} \qquad Warmest \blank{2.0cm}
      \item Is $66^\circ$F acceptable? \blank{1.8cm} \; Which word in the sentence tells you?
            \blank{3.2cm}
    \end{enumerate}

  \item \textbf{Multiple choice.} Which is the solution set of $|x + 2| > 3$? Justify in one
        sentence.
    \begin{enumerate}[label=\Alph*., leftmargin=2.2em, itemsep=1pt]
      \item $-5 < x < 1$
      \item $x < -5$ \; or \; $x > 1$
      \item $x > 1$
      \item all real numbers
    \end{enumerate}
    \answerspace{1.2cm}{}
\end{enumerate}
\end{notesbox}

\boxguard[10]
\begin{spiralbox}
\textbf{Coming next --- Lesson 1.3: Absolute Value Functions \& Transformations.} Today you
\emph{solved} $|ax+b| = c$; next you will \emph{graph} $y = |ax+b|$ --- a \textbf{V-shape}. The two
solutions you found for an equation are exactly where that V crosses the horizontal line $y = c$,
which previews the vertex, axis of symmetry, and minimum you will study formally in Unit~2.
\end{spiralbox}

\end{document}
